%=====================================================================
\chapter{Bayesian Learning and Naive Bayes}\label{ch:bayes}
%=====================================================================
\locator{3}{Part III --- Supervised Learning Algorithms}

\begin{outcomes}
By the end of this chapter you will be able to:
\begin{tight}
  \item \textbf{Apply} Bayes' theorem to hypotheses, and \textbf{find} the MAP
        and maximum likelihood hypotheses.
  \item \textbf{Explain} the Bayes optimal classifier and why it can disagree
        with the single most probable hypothesis.
  \item \textbf{Classify} an example with naive Bayes by hand, and
        \textbf{state} its conditional independence assumption.
  \item \textbf{Apply} Laplace smoothing, and \textbf{use} Gaussian naive Bayes
        for continuous features.
  \item \textbf{Build} a multinomial naive Bayes text classifier working in
        log-probabilities.
\end{tight}
\end{outcomes}

\begin{prereq}
\chref{math} (conditional probability, Bayes' theorem, the Gaussian, maximum
likelihood) and \chref{trees} (the sprint data, used again here).
\end{prereq}

\begin{keyterms}
prior $\bullet$ likelihood $\bullet$ posterior $\bullet$ evidence $\bullet$ MAP
hypothesis $\bullet$ maximum likelihood hypothesis $\bullet$ Bayes optimal
classifier $\bullet$ conditional independence $\bullet$ naive Bayes $\bullet$
zero-frequency problem $\bullet$ Laplace smoothing $\bullet$ $m$-estimate
$\bullet$ Gaussian naive Bayes $\bullet$ bag of words $\bullet$ multinomial naive
Bayes $\bullet$ log-probabilities
\end{keyterms}

\section{Learning as Updating Belief}

The learners so far return one hypothesis and no statement of how sure they are.
Bayesian learning treats every hypothesis as uncertain: it starts from a
\term{prior} belief about which hypotheses are plausible, observes data, and
computes a \term{posterior} belief. Two things follow that the earlier chapters
could not offer. Prior knowledge can be put in explicitly, and every prediction
comes with a probability. And even where Bayesian methods are too expensive to
use directly, they give the standard against which other methods are judged.

\section{Bayes' Theorem for Hypotheses}

\begin{definitionbox}[Posterior, MAP and ML Hypotheses]
For a hypothesis $h$ and training data $D$:
\[
P(h \mid D) = \frac{P(D \mid h)\,P(h)}{P(D)}
\]
$P(h)$ is the \term{prior}, $P(D \mid h)$ the \term{likelihood}, $P(D)$ the
\term{evidence} and $P(h \mid D)$ the \term{posterior}. The \term{maximum a
posteriori} (MAP) hypothesis maximises the posterior; since $P(D)$ is the same for
every $h$, it can be dropped:
\[
h_{MAP} = \arg\max_{h \in H} P(D \mid h)\,P(h)
\]
If every hypothesis is equally likely a priori, this reduces to the \term{maximum
likelihood} hypothesis $h_{ML} = \arg\max_h P(D \mid h)$.
\end{definitionbox}

\begin{worked}[A Flagged Code Change]
Two hypotheses: a code change contains a critical defect, or it is clean. 0.8\% of
changes contain one. A static-analysis tool flags 98\% of the changes that do, and
3\% of those that do not. The tool flags a change.

\smallskip
\begin{tight}
  \item $P(\text{flag} \mid \text{defect})\,P(\text{defect}) = 0.98 \times 0.008 = 0.0078$
  \item $P(\text{flag} \mid \text{clean})\,P(\text{clean}) = 0.03 \times 0.992 = 0.0298$
\end{tight}
\textbf{$h_{MAP} =$ clean}, because $0.0298 > 0.0078$. Normalising:
\[
P(\text{defect} \mid \text{flag}) = \frac{0.0078}{0.0078 + 0.0298} = \mathbf{0.21}
\]
A flag raises the probability from 0.8\% to 21\% --- a twenty-six-fold increase
--- and it is still more likely than not that the change is clean. The maximum
likelihood hypothesis, which ignores the prior, would say ``defect''
($0.98 > 0.03$). This is the base-rate effect of \chref{math}, now read as a choice
between hypotheses.
\end{worked}

\dsfig{%
\begin{tikzpicture}
\begin{groupplot}[
  group style={group size=3 by 1, horizontal sep=1.0cm},
  width=4.9cm, height=4.4cm, ybar, /pgf/bar width=0.55cm, axis lines=left,
  ymin=0, symbolic x coords={defect,clean}, xtick=data,
  tick label style={font=\tiny}, title style={font=\scriptsize\bfseries},
  nodes near coords, nodes near coords style={font=\tiny},
  enlarge x limits=0.5]
\nextgroupplot[title={prior $P(h)$}, ymax=1.2,
  point meta=explicit symbolic]
\addplot[fill=primary!45, draw=primary] coordinates
  {(defect,0.008)[0.008] (clean,0.992)[0.992]};
\nextgroupplot[title={likelihood $P(\text{flag} \mid h)$}, ymax=1.2,
  point meta=explicit symbolic]
\addplot[fill=goldhl!55, draw=goldhl] coordinates
  {(defect,0.98)[0.98] (clean,0.03)[0.03]};
\nextgroupplot[title={posterior $P(h \mid \text{flag})$}, ymax=1.0,
  point meta=explicit symbolic]
\addplot[fill=greenhl!50, draw=greenhl!70!black] coordinates
  {(defect,0.21)[0.21] (clean,0.79)[0.79]};
\end{groupplot}
\end{tikzpicture}}%
{Prior times likelihood, normalised, gives the posterior. A flag that strongly favours
``defect'' is not enough to overturn a prior that strongly favours ``clean''.}{bayes-update}

\begin{conceptbox}[Least Squares Is a MAP Hypothesis in Disguise]
If the training labels are the true function plus Gaussian noise and all
hypotheses are equally likely a priori, the maximum likelihood hypothesis is the
one that minimises the sum of squared errors --- exactly the least-squares fit of
\chref{linreg}. Put a Gaussian prior on the weights as well, favouring small
ones, and the MAP hypothesis is the ridge regression of \chref{generalization}.
Regularisation is a prior.
\end{conceptbox}

\section{The Bayes Optimal Classifier}

The MAP hypothesis answers ``which hypothesis is most probable?''. That is not the
same question as ``which classification is most probable?''. The \term{Bayes optimal
classifier} answers the second by letting every hypothesis vote, weighted by its
posterior:
\[
v_{\text{opt}} = \arg\max_{v} \sum_{h \in H} P(v \mid h)\,P(h \mid D)
\]
No other method using the same hypothesis space and prior can do better on
average --- hence ``optimal''. It is rarely computable, because $H$ is too large to
sum over, but it is the benchmark every approximation aims at.

\begin{worked}[When the Most Probable Hypothesis Is Outvoted]
Three hypotheses remain with posteriors $P(h_1 \mid D) = 0.4$, $P(h_2 \mid D) = 0.3$
and $P(h_3 \mid D) = 0.3$. A new instance is classified positive by $h_1$ and negative
by $h_2$ and $h_3$.
\begin{tight}
  \item \textbf{MAP:} $h_1$ is the most probable hypothesis, so ``positive''.
  \item \textbf{Bayes optimal:} $P(+) = 0.4$ and $P(-) = 0.3 + 0.3 = 0.6$, so
        \textbf{``negative''}.
\end{tight}
The single best hypothesis is outvoted by the combined weight of the others. The
same idea --- let many hypotheses vote --- is the basis of the ensembles of
\chref{ensembles}.
\end{worked}

\section{Naive Bayes}

For classification with attributes $a_1, \dots, a_d$, the most probable class is
\[
v_{MAP} = \arg\max_{v} P(v)\,P(a_1, a_2, \dots, a_d \mid v)
\]
The prior $P(v)$ is easy to estimate: count. The joint likelihood is not: there are
as many combinations of attribute values as there are possible examples, and most
never appear in the training data. Naive Bayes removes the problem with one bold
assumption.

\begin{definitionbox}[The Naive Bayes Classifier]
Assume the attributes are \term{conditionally independent} given the class:
$P(a_1, \dots, a_d \mid v) = \prod_j P(a_j \mid v)$. Then
\[
v_{NB} = \arg\max_{v}\; P(v)\prod_{j=1}^{d} P(a_j \mid v)
\]
Each $P(a_j \mid v)$ is estimated from counts: the fraction of class-$v$ training
examples that have value $a_j$. Training is a single pass of counting.
\end{definitionbox}

\dsfig{%
\begin{tikzpicture}[font=\scriptsize,
  cls/.style={circle, draw=primary, fill=primary!14, text=primary!80!black,
              font=\scriptsize\bfseries, minimum size=1.3cm, align=center},
  att/.style={circle, draw=secondary, fill=secondary!8, text=secondary!70!black,
              font=\scriptsize, minimum size=1.15cm, align=center},
  a/.style={-{Stealth[length=2mm]}, gray!60, line width=0.9pt}]
\node[cls] (c) at (0,0) {Goal\\met?};
\node[att] (o) at (-3.3,-2.1) {Reqs.};
\node[att] (t) at (-1.1,-2.1) {Pressure};
\node[att] (h) at (1.1,-2.1) {Load};
\node[att] (w) at (3.3,-2.1) {Deps.};
\foreach \n in {o,t,h,w}{\draw[a] (c) -- (\n);}
\node[font=\tiny, text=gray!55!black, align=center] at (0,-3.2)
     {no arrows between the attributes: given the class, they are assumed independent};
\node[anchor=west, align=left, text width=5.8cm, font=\scriptsize, text=gray!55!black]
     at (4.6,-0.9)
     {\textbf{What is estimated.} One prior $P(v)$ per class, and one table
      $P(a_j \mid v)$ per attribute and class. For the sprint data that is 2 priors
      and $2 \times (3 + 3 + 2 + 2) = 20$ conditional probabilities --- instead of a
      joint distribution over all 36 combinations of attribute values for each
      class.};
\end{tikzpicture}}%
{The naive Bayes assumption drawn as a network: the class causes each attribute, and
the attributes do not influence each other once the class is known.}{nb-graph}

\begin{worked}[Naive Bayes on the Sprint Data]
Will a new sprint with $\langle$Changing requirements, Low pressure, High load,
Many dependencies$\rangle$ meet its goal? Use the fourteen sprints of \chref{trees}.

\smallskip
\textbf{Priors:} $P(\text{Yes}) = 9/14 = 0.643$; $P(\text{No}) = 5/14 = 0.357$.

\textbf{Conditionals}, counted within each class:
\begin{tight}
  \item Yes: $P(\text{Changing}\mid\text{Yes}) = 2/9$, $P(\text{Low}\mid\text{Yes}) = 3/9$,
        $P(\text{High load}\mid\text{Yes}) = 3/9$, $P(\text{Many}\mid\text{Yes}) = 3/9$.
  \item No: $P(\text{Changing}\mid\text{No}) = 3/5$, $P(\text{Low}\mid\text{No}) = 1/5$,
        $P(\text{High load}\mid\text{No}) = 4/5$, $P(\text{Many}\mid\text{No}) = 3/5$.
\end{tight}
\textbf{Scores:}
\[
\text{Yes: } 0.643 \times \tfrac{2}{9} \times \tfrac{3}{9} \times \tfrac{3}{9} \times \tfrac{3}{9} = 0.0053
\qquad
\text{No: } 0.357 \times \tfrac{3}{5} \times \tfrac{1}{5} \times \tfrac{4}{5} \times \tfrac{3}{5} = 0.0206
\]
\textbf{Prediction: No.} Normalised, $P(\text{No} \mid x) = 0.0206/(0.0206 + 0.0053)
= \mathbf{0.795}$. (The decision tree of \chref{trees} agrees: changing
requirements and a high load means the goal is missed.)
\end{worked}

\section{The Zero-Frequency Problem and Laplace Smoothing}

Classify $\langle$Frozen, High, High, Many$\rangle$. No training sprint that missed
its goal had Frozen requirements, so $P(\text{Frozen} \mid \text{No}) = 0/5 = 0$,
and the whole product for No is zero --- however strongly a high load and many
dependencies point towards No. One unseen attribute value silences all the others.

\begin{definitionbox}[Laplace Smoothing and the $m$-Estimate]
Add one imaginary example of every value:
\[
P(a_j = v \mid c) = \frac{n_{v,c} + 1}{n_c + k_j}
\]
where $n_{v,c}$ counts class-$c$ examples with that value, $n_c$ counts class-$c$
examples, and $k_j$ is the number of values attribute $j$ can take.
\end{definitionbox}

Mitchell's more general \term{$m$-estimate}, $(n_{v,c} + m\,p)/(n_c + m)$, adds $m$
imaginary examples distributed according to a prior estimate $p$; Laplace smoothing
is the case $m = k_j$, $p = 1/k_j$.

\begin{worked}[Smoothing Rescues the Frozen Sprint]
With Laplace smoothing (Requirements and Pressure have 3 values, Load and
Dependencies 2):
\begin{tight}
  \item Yes: $0.643 \times \tfrac{4+1}{9+3} \times \tfrac{2+1}{9+3} \times
        \tfrac{3+1}{9+2} \times \tfrac{3+1}{9+2} = 0.643 \times 0.417 \times 0.250
        \times 0.364 \times 0.364 = 0.0089$
  \item No: $0.357 \times \tfrac{0+1}{5+3} \times \tfrac{2+1}{5+3} \times
        \tfrac{4+1}{5+2} \times \tfrac{3+1}{5+2} = 0.357 \times 0.125 \times 0.375
        \times 0.714 \times 0.571 = 0.0068$
\end{tight}
$P(\text{Yes} \mid x) = 0.0089/(0.0089 + 0.0068) = \mathbf{0.56}$. Without smoothing the
answer was ``Yes with certainty''; with it, a narrow Yes that reflects the conflict
between Frozen requirements (always Yes in training) and a high load with many
dependencies (usually No).
\end{worked}

\section{Continuous Features: Gaussian Naive Bayes}

For a numeric attribute, counting values is useless --- every value is new. Instead
model each attribute within each class as a Gaussian, with the class's own mean
$\mu_{j,c}$ and standard deviation $\sigma_{j,c}$ estimated from the training data,
and use the density as the likelihood:
\[
P(a_j \mid c) \;\propto\; \frac{1}{\sqrt{2\pi}\,\sigma_{j,c}}
\exp\!\left(-\frac{(a_j - \mu_{j,c})^2}{2\sigma_{j,c}^2}\right)
\]

\begin{worked}[One Feature, Two Classes, a Strong Prior]
Modules that pass QA see 500 lines of code changed per week on average (sd 100);
modules that go on to fail QA see 900 (sd 150); 5\% of modules fail. This week a
module had 750 lines changed.
\begin{tight}
  \item Pass: density $\dfrac{1}{100\sqrt{2\pi}}e^{-(250)^2/(2 \cdot 100^2)}
        = 0.00399 \times e^{-3.125} = 0.000175$
  \item Fail: density $\dfrac{1}{150\sqrt{2\pi}}e^{-(150)^2/(2 \cdot 150^2)}
        = 0.00266 \times e^{-0.5} = 0.00161$
\end{tight}
The likelihood favours failure by $0.00161/0.000175 = 9.2$ to 1. Now the priors:
\[
P(\text{fail} \mid 750) = \frac{0.05 \times 0.00161}{0.05 \times 0.00161 + 0.95 \times 0.000175}
= \frac{0.0000807}{0.000247} = \mathbf{0.33}
\]
Nine-to-one evidence against nineteen-to-one prior odds leaves the module more
likely to pass. \dsref{gaussian-nb} shows the whole picture: the boundary is at
772 lines.
\end{worked}

\dsfig{%
\begin{tikzpicture}
\begin{groupplot}[
  group style={group size=2 by 1, horizontal sep=1.4cm},
  width=6.9cm, height=4.8cm, axis lines=left,
  tick label style={font=\tiny}, label style={font=\tiny},
  title style={font=\scriptsize\bfseries},
  xlabel={lines changed per week}, xmin=150, xmax=1450]
\nextgroupplot[title={Likelihoods $P(x \mid \text{class})$}, ylabel={density},
  ymin=0, ymax=0.0046, ytick=\empty,
  legend style={font=\tiny, draw=gray!40, at={(0.98,0.97)}, anchor=north east}]
\addplot[primary, line width=1.3pt, domain=150:1450, samples=150]
  {exp(-(x-500)^2/(2*100^2))/(100*sqrt(2*pi))};
\addlegendentry{passes QA}
\addplot[accent, line width=1.3pt, domain=150:1450, samples=150]
  {exp(-(x-900)^2/(2*150^2))/(150*sqrt(2*pi))};
\addlegendentry{fails QA}
\draw[gray!55, dashed] (axis cs:750,0) -- (axis cs:750,0.0030);
\node[font=\tiny, gray!55!black, anchor=south] at (axis cs:750,0.0030) {750};
\nextgroupplot[title={Posterior $P(\text{fail} \mid x)$, prior 5\%},
  ylabel={probability}, ymin=0, ymax=1.05]
\addplot[accent, line width=1.3pt, domain=150:1450, samples=200]
  {0.05*exp(-(x-900)^2/(2*150^2))/150 /
   (0.05*exp(-(x-900)^2/(2*150^2))/150 + 0.95*exp(-(x-500)^2/(2*100^2))/100)};
\addplot[gray!55, dashed] coordinates {(150,0.5)(1450,0.5)};
\draw[greenhl!70!black, line width=0.9pt] (axis cs:772,0) -- (axis cs:772,0.5);
\node[font=\tiny, greenhl!45!black, anchor=north west, align=left] at (axis cs:790,0.44)
     {boundary at 772 lines:\\beyond it, fail wins};
\addplot[only marks, mark=*, mark size=1.8pt, primary] coordinates {(750,0.326)};
\node[font=\tiny, primary, anchor=east] at (axis cs:730,0.33) {750 lines: 0.33};
\end{groupplot}
\end{tikzpicture}}%
{Left: the two class-conditional densities; at 750 lines the fail density is nine times
higher. Right: after weighting by the 5\% prior, the posterior crosses 0.5 only at 772
lines. The boundary sits beyond the midpoint of the means because the rare class needs
more evidence.}{gaussian-nb}

\section{Text Classification}

Naive Bayes's most famous success is text: routing tickets, labelling commit
messages, spotting duplicate bug reports. A document is represented as a
\term{bag of words} --- the count of each vocabulary word, order ignored --- and
\term{multinomial naive Bayes} treats each word occurrence as an independent draw
from a class-specific distribution over the vocabulary:
\begin{gather*}
P(w \mid c) = \frac{\text{count}(w, c) + 1}{\text{total words in } c + |V|}, \\
v_{NB} = \arg\max_c\Big[\log P(c) + \sum_{\text{words } w \text{ in doc}} \log P(w \mid c)\Big]
\end{gather*}
Working in logarithms turns the product into a sum; a product of a few hundred
probabilities would otherwise underflow to zero on a computer.

\begin{worked}[A Four-Word Ticket Router]
Training counts over a vocabulary of $|V| = 4$ words, from bug reports and feature
requests; both classes contain 40 words and are equally likely:

\smallskip
\begin{center}
\begin{tabular}{@{}lrrrr@{}}
\toprule
 & \textbf{crash} & \textbf{error} & \textbf{add} & \textbf{option} \\
\midrule
bug & 20 & 15 & 2 & 3 \\
feature & 2 & 1 & 18 & 19 \\
\bottomrule
\end{tabular}
\end{center}

\smallskip
With smoothing, $P(\text{crash}\mid\text{bug}) = 21/44$, $P(\text{add}\mid\text{bug}) = 3/44$,
$P(\text{crash}\mid\text{feature}) = 3/44$, $P(\text{add}\mid\text{feature}) = 19/44$.

\textbf{Ticket:} ``crash add crash''.
\begin{tight}
  \item bug: $\log\tfrac{1}{2} + 2\log\tfrac{21}{44} + \log\tfrac{3}{44}
        = -0.693 - 1.479 - 2.686 = -4.858$
  \item feature: $\log\tfrac{1}{2} + 2\log\tfrac{3}{44} + \log\tfrac{19}{44}
        = -0.693 - 5.371 - 0.840 = -6.904$
\end{tight}
\textbf{Bug} wins by $2.05$ in log-probability, a likelihood ratio of
$e^{2.05} \approx 7.7$, so $P(\text{bug}) = 7.7/8.7 = \mathbf{0.885}$. The two
``crash''es outweigh the one ``add'' --- which also shows the router's weakness: a
feature request that happens to mention crashes is sent to the bug queue.
\end{worked}

\begin{alertbox}[Why Naive Bayes Works When Its Assumption Is False]
Words are obviously not independent --- ``stack'' and ``trace'' travel together.
Naive Bayes therefore double-counts correlated evidence and its probabilities are
badly \emph{over-confident}. But classification needs only the right \emph{ranking}
of classes, and the double-counting usually inflates the right class as much as
it helps it. Use naive Bayes's decisions freely; do not trust its probabilities
without calibration. Models that represent the dependencies explicitly --- Bayesian
networks --- are the next step, and are covered in Mitchell's Chapter~6.
\end{alertbox}

\begin{pitfall}
\begin{tight}
  \item \textbf{Forgetting to smooth.} One zero count vetoes the class.
  \item \textbf{Multiplying hundreds of probabilities.} Add log-probabilities.
  \item \textbf{Confusing MAP with ML.} Maximum likelihood ignores the prior, and
        with rare classes the prior decides.
  \item \textbf{Reporting naive Bayes probabilities as calibrated risks.} They are
        too extreme; the ranking is what is reliable.
  \item \textbf{Feeding Gaussian NB heavily skewed features.} Transform them (log)
        first, or the Gaussian fits badly.
\end{tight}
\end{pitfall}

%---------------------------------------------------------------------
\begin{lab}[Naive Bayes by Counting, Then by Library]
\begin{lstlisting}[language=Python]
import numpy as np
from sklearn.naive_bayes import GaussianNB, MultinomialNB
from sklearn.feature_extraction.text import CountVectorizer

# --- 1. naive Bayes on the sprint data, by counting ----------------------
sprints = [("Changing High High Few", "No"),
           ("Changing High High Many", "No"),
           ("Frozen High High Few", "Yes"),
           ("Vague Medium High Few", "Yes"),
           ("Vague Low Normal Few", "Yes"), ("Vague Low Normal Many", "No"),
           ("Frozen Low Normal Many", "Yes"),
           ("Changing Medium High Few", "No"),
           ("Changing Low Normal Few", "Yes"),
           ("Vague Medium Normal Few", "Yes"),
           ("Changing Medium Normal Many", "Yes"),
           ("Frozen Medium High Many", "Yes"),
           ("Frozen High Normal Few", "Yes"),
           ("Vague Medium High Many", "No")]
X = [s.split() for s, _ in sprints]
y = [c for _, c in sprints]
n_values = [3, 3, 2, 2]                  # values per attribute


def naive_bayes(x, laplace=0):
    scores = {}
    for c in ("Yes", "No"):
        rows = [r for r, lab in zip(X, y) if lab == c]
        p = len(rows) / len(X)                            # prior P(c)
        for j, v in enumerate(x):                         # P(x_j | c)
            count = sum(r[j] == v for r in rows)
            p *= (count + laplace) / (len(rows) + laplace * n_values[j])
        scores[c] = p
    total = sum(scores.values())
    return {c: round(s / total, 3) for c, s in scores.items()}, scores


post, raw = naive_bayes(["Changing", "Low", "High", "Many"])
print("unnormalised:", {c: round(s, 5) for c, s in raw.items()})
print("posterior:   ", post)
frozen = ["Frozen", "High", "High", "Many"]
print("Frozen sprint, no smoothing:", naive_bayes(frozen)[1])
print("Frozen sprint, Laplace     :", naive_bayes(frozen, laplace=1)[0])

# --- 2. Gaussian naive Bayes on the code-churn example -------------------
rng = np.random.default_rng(0)
churn = np.r_[rng.normal(500, 100, 1900),       # modules that pass QA
              rng.normal(900, 150, 100)]         # the 5% that fail
fails = np.r_[np.zeros(1900), np.ones(100)]
gnb = GaussianNB().fit(churn.reshape(-1, 1), fails)
print("learned means:", gnb.theta_.ravel().round(0))
print("P(fail | 750 lines):", gnb.predict_proba([[750]])[0, 1].round(2))
grid = np.arange(600, 1000).reshape(-1, 1)
print("boundary near", int(grid[gnb.predict(grid) == 1][0, 0]), "lines")

# --- 3. a ticket router from twelve tickets ------------------------------
tickets = ["app crashes when saving", "error on the export page",
           "crash after the update", "null pointer error in upload",
           "report page crashes on load", "error message on checkout",
           "add an option to export pdf", "please add a dark mode",
           "option to filter by date", "add support for csv import",
           "add a bulk edit option", "please add weekly summary"]
labels = [1] * 6 + [0] * 6                         # 1 = bug
vec = CountVectorizer()
nb = MultinomialNB(alpha=1.0).fit(vec.fit_transform(tickets), labels)
for t in ["crash after export", "add export option", "crash option crash"]:
    p = nb.predict_proba(vec.transform([t]))[0, 1]
    print(f"{t!r:22s} P(bug) = {p:.3f}")
\end{lstlisting}
Part 1 reproduces both worked examples. Part 2 samples modules from the two
Gaussians of the worked example and lets \texttt{GaussianNB} estimate them; its
answer is close to the 0.33 and 772 lines computed by hand. In part 3,
\texttt{alpha} is the smoothing constant: set it to \texttt{1e-9} and classify
``crash option crash'' again.
\end{lab}

%---------------------------------------------------------------------
\begin{checkpoint}
\begin{tightnum}
  \item 1\% of releases would break production. A pre-release check flags 90\% of
        those and 5\% of good releases. Give $h_{MAP}$ and $h_{ML}$ for a release
        the check has flagged.
  \item What exactly does naive Bayes assume, and what goes wrong in the
        probabilities when the assumption fails?
  \item Why must naive Bayes for text work with log-probabilities?
\end{tightnum}
\end{checkpoint}

%---------------------------------------------------------------------
\begin{chaptersummary}
\begin{tight}
  \item Bayesian learning updates a prior into a posterior with the likelihood;
        $h_{MAP}$ maximises $P(D\mid h)P(h)$ and equals $h_{ML}$ under a uniform prior.
  \item Least squares is maximum likelihood under Gaussian noise; ridge is MAP
        with a Gaussian prior on the weights.
  \item The Bayes optimal classifier weights every hypothesis's vote by its
        posterior and can overrule the MAP hypothesis.
  \item Naive Bayes assumes attributes are independent given the class, which
        reduces training to counting; on the sprint data it predicts No with 0.795.
  \item Smooth the counts (Laplace, $m$-estimate) to avoid zero vetoes; use
        Gaussian likelihoods for continuous features.
  \item For text, use a bag of words and log-probabilities. Trust the ranking, not
        the probabilities.
\end{tight}
\end{chaptersummary}

\begin{reviewq}
  \item \textbf{(MCQ)} With a uniform prior over hypotheses, $h_{MAP}$ equals:
        (a)~$h_{ML}$ (b)~the Bayes optimal classifier (c)~the most specific
        hypothesis (d)~none of these.
  \item \textbf{(MCQ)} Laplace smoothing is used to: (a)~speed up training
        (b)~prevent zero probabilities for unseen values (c)~normalise the
        posterior (d)~remove correlated features.
  \item \textbf{(MCQ)} In Gaussian naive Bayes, each class--feature pair is
        described by: (a)~a count table (b)~a mean and a variance (c)~a
        threshold (d)~a weight.
  \item \textbf{(Short)} Explain the difference between the MAP hypothesis and the
        Bayes optimal classification, with an example.
  \item \textbf{(Short)} Why can naive Bayes classify well even when its
        probabilities are poor?
  \item \textbf{(Short)} Describe how a user could word a feature request so that
        a naive Bayes router sends it to the faster bug queue.
  \item \textbf{(Applied)} Using the sprint data and Laplace smoothing, classify
        $\langle$Vague, High, Normal, Few$\rangle$. Show every probability.
  \item \textbf{(Applied)} For the Gaussian example, find the weekly churn at which
        the two prior-weighted densities are equal, to the nearest 10 lines, by
        trying values. Explain why it is not 700, the midpoint of the means.
\end{reviewq}
